% Generated. Do not edit.
% Projection scriptural-date-location/section, version 1: the Scriptural Date and
% Location page, generated whole from research/chronology.toml, the Scripture
% chronology corpus and the source library by tools/tpt document-types refresh.
% document: liturgy/roman-rite/postconciliar/roman-missal-third-edition-en-us-2011/propers/temporal/pc-s55-twenty-ninth-sunday-in-ordinary-time-year-a
\runninghead{Scriptural Date and Location}
\section{Scriptural Date and Location}
\label{sec:date-location}

{\small Each passage the Mass reads or sings from Scripture, in canonical order. The Context column names what is dated --- the event the passage narrates, a setting it remembers, the traditional attribution of its author, or the composition of its book --- and where it belongs; the Date column gives the answer of the Scripture chronology corpus for that subject, with every alternative it records. Texts of the Mass that are not drawn from Scripture have no row. The evidence for each date follows the table.\par}
\medskip
% Generated. Do not edit.
% Deterministic projection of research/chronology.toml.
% schema: 5
% document: liturgy/roman-rite/postconciliar/roman-missal-third-edition-en-us-2011/propers/temporal/pc-s55-twenty-ninth-sunday-in-ordinary-time-year-a
% generated-by: tools/tpt proper-chronology annotations
\newcommand{\chronologyannotationclaim}[6]{#6}
\newcommand{\chronologyannotationreach}[3]{}
\newcommand{\chronologyannotationgroup}[3]{#3}
\newcommand{\chronologyannotation}[1]{%
  \ifcsname triptychchronologyannotation@#1\endcsname
    \csname triptychchronologyannotation@#1\endcsname
  \else
    \PackageError{triptych}{No chronology annotation for #1}{}%
  \fi}
% appointment-dependency: guidance/liturgy/postconciliar-propers-registry.md sha256=376860a7c37a4c75bb6d880cd1b260183be1877a07e3cd01660a9e2854e2ea72
% appointment-dependency: src/gpt/liturgy/roman-rite/postconciliar/roman-missal-third-edition-en-us-2011/propers/registry/README.md sha256=4e2d92326dece103e1475e70ffb5e447815f9b12d697926e48a54850a1641208
% appointment-dependency: src/gpt/liturgy/roman-rite/postconciliar/roman-missal-third-edition-en-us-2011/propers/registry/formula-dispositions.md sha256=de58754f5a8ab61dffa3f4a6a479c843fe3763b6a5b37099282960de71d93989
% appointment-dependency: src/gpt/liturgy/roman-rite/postconciliar/roman-missal-third-edition-en-us-2011/propers/temporal/pc-s55-twenty-ninth-sunday-in-ordinary-time-year-a/instance/manifest.md sha256=a33c95bbbb2f9b954e8f879cdcb159ed91b18220a86e6917c9a6092c1e6d32fd
% appointment-dependency: src/gpt/liturgy/roman-rite/postconciliar/roman-missal-third-edition-en-us-2011/propers/temporal/pc-s55-twenty-ninth-sunday-in-ordinary-time-year-a/propers/verified.md sha256=4e5b677a3ceeffc388e41c1d9e2e31b690200ab9ba63365e88e975229db7de2a
% appointment-dependency: src/gpt/liturgy/roman-rite/postconciliar/roman-missal-third-edition-en-us-2011/propers/temporal/pc-s55-twenty-ninth-sunday-in-ordinary-time-year-a/research/chronology-inputs.toml sha256=0b689b16340810b8d2a46dcf681626b0b7ba23be25e11a5f8fcf033b1d33dce4
% appointment-dependency: src/gpt/liturgy/roman-rite/postconciliar/roman-missal-third-edition-en-us-2011/propers/temporal/shared/ordinary-time/weeks/29/propers/verified.md sha256=1e6ac1e5369a1ab134cc2969d1caf70f7dead1c18610987f6b89d80c11d089b1
% appointment-note: entrance: Cf. Psalm 16:6,8; adaptation; vulgate numbering; canonical verse queries, not appointed wording.
% appointment-note: collect: Composed Missal oration without an appointed biblical citation; theological affinities do not create appointments.
% appointment-note: first-reading: Isaiah 45:1,4-6; appointed; vulgate numbering; canonical verse queries, not appointed wording.
% appointment-note: psalm: Psalm 96:1,3,4-5,7-8,9,10a,10c; appointed; hebrew numbering; whole-verse query envelope, not subverse chronology or appointed wording.
% appointment-note: psalm: Psalm 96:7b; response; hebrew numbering; whole-verse query envelope, not subverse chronology or appointed wording.
% appointment-note: second-reading: 1 Thessalonians 1:1-5b; appointed; vulgate numbering; whole-verse query envelope, not subverse chronology or appointed wording.
% appointment-note: acclamation: Philippians 2:15d,16a; appointed; vulgate numbering; whole-verse query envelope, not subverse chronology or appointed wording.
% appointment-note: gospel: Matthew 22:15-21; appointed; vulgate numbering; canonical verse queries, not appointed wording.
% appointment-note: offerings: Composed Missal oration without an appointed biblical citation; theological affinities do not create appointments.
% appointment-note: communion-psalm: Cf. Psalm 32:18-19; adaptation; vulgate numbering; canonical verse queries, not appointed wording.
% appointment-note: communion-mark: Mark 10:45; identified-basis; vulgate numbering; canonical verse queries, not appointed wording.
% appointment-note: after-communion: Composed Missal oration without an appointed biblical citation; theological affinities do not create appointments.
% entrance: Entrance antiphon
\expandafter\def\csname triptychchronologyannotation@entrance\endcsname{%
  \chronologyannotationgroup{historical-setting}{research-pending}{\textbf{Historical setting} -- Occasion date unresolved.}%
  \space \chronologyannotationgroup{composition}{preferred}{\textbf{Composition}: \chronologyannotationreach{Ps.16.6}{Ps}{true}\chronologyannotationreach{Ps.16.8}{Ps}{true}\chronologyannotationclaim{critical.psalms.latest-composition-boundary}{composition}{catholic-critical-v1}{preferred}{before the Maccabean period, around 165 B.C.; no individual psalm can be dated securely}{Before c. 165 B.C.}}%
}
% first-reading: First reading
\expandafter\def\csname triptychchronologyannotation@first-reading\endcsname{%
  \chronologyannotationgroup{traditional-attribution}{preferred}{\textbf{Traditional attribution}: \chronologyannotationreach{Is.45.1}{Is}{true}\chronologyannotationreach{Is.45.4}{Is}{true}\chronologyannotationreach{Is.45.5}{Is}{true}\chronologyannotationreach{Is.45.6}{Is}{true}\chronologyannotationclaim{israel.prophets.isaias-traditional-era}{traditional-attribution}{catholic-traditional-v1}{preferred}{from the closing year of Ozias, King of Juda (740 B.C.); no prophecies appear to be later than 701 B.C.}{Isaias (ministry in Souvay's traditional account), B.C. 740--701}.}%
  \space \chronologyannotationgroup{prophecy-given}{preferred}{\textbf{Prophecy given}: \chronologyannotationreach{Is.45.1}{Is.40 Is.41 Is.42 Is.43 Is.44 Is.45 Is.46 Is.47 Is.48 Is.49 Is.50 Is.51 Is.52 Is.53 Is.54 Is.55}{true}\chronologyannotationreach{Is.45.4}{Is.40 Is.41 Is.42 Is.43 Is.44 Is.45 Is.46 Is.47 Is.48 Is.49 Is.50 Is.51 Is.52 Is.53 Is.54 Is.55}{true}\chronologyannotationreach{Is.45.5}{Is.40 Is.41 Is.42 Is.43 Is.44 Is.45 Is.46 Is.47 Is.48 Is.49 Is.50 Is.51 Is.52 Is.53 Is.54 Is.55}{true}\chronologyannotationreach{Is.45.6}{Is.40 Is.41 Is.42 Is.43 Is.44 Is.45 Is.46 Is.47 Is.48 Is.49 Is.50 Is.51 Is.52 Is.53 Is.54 Is.55}{true}\chronologyannotationclaim{israel.exile.second-isaias-oracles}{prophecy-given}{catholic-critical-v1}{preferred}{toward the end of the Babylonian exile}{The prophetic activity behind Isaiah 40-55 in the NABRE introduction, toward the end of the Babylonian exile}.}%
}
% psalm: Responsorial Psalm
\expandafter\def\csname triptychchronologyannotation@psalm\endcsname{%
  \chronologyannotationgroup{superscription-setting}{preferred}{\textbf{Superscription setting}: \chronologyannotationreach{Ps.95.1}{Ps.95}{true}\chronologyannotationreach{Ps.95.10}{Ps.95}{true}\chronologyannotationreach{Ps.95.3}{Ps.95}{true}\chronologyannotationreach{Ps.95.4}{Ps.95}{true}\chronologyannotationreach{Ps.95.5}{Ps.95}{true}\chronologyannotationreach{Ps.95.7}{Ps.95}{true}\chronologyannotationreach{Ps.95.8}{Ps.95}{true}\chronologyannotationreach{Ps.95.9}{Ps.95}{true}\chronologyannotationclaim{israel.restoration.temple-raised-by-zorobabel}{superscription-setting}{catholic-traditional-v1}{preferred}{(537 B.C.)}{The raising of the second Temple by Zorobabel, B.C. 537}.}%
  \space \chronologyannotationgroup{composition}{preferred}{\textbf{Composition}: \chronologyannotationreach{Ps.95.1}{Ps}{true}\chronologyannotationreach{Ps.95.10}{Ps}{true}\chronologyannotationreach{Ps.95.3}{Ps}{true}\chronologyannotationreach{Ps.95.4}{Ps}{true}\chronologyannotationreach{Ps.95.5}{Ps}{true}\chronologyannotationreach{Ps.95.7}{Ps}{true}\chronologyannotationreach{Ps.95.8}{Ps}{true}\chronologyannotationreach{Ps.95.9}{Ps}{true}\chronologyannotationclaim{critical.psalms.latest-composition-boundary}{composition}{catholic-critical-v1}{preferred}{before the Maccabean period, around 165 B.C.; no individual psalm can be dated securely}{Before c. 165 B.C.}}%
  \space \chronologyannotationgroup{prophetic-referent}{disputed}{\textbf{Prophetic referent} -- disputed: \chronologyannotationreach{Ps.95.1}{Ps.67 Ps.95 Ps.96 Ps.97}{true}\chronologyannotationreach{Ps.95.10}{Ps.67 Ps.95 Ps.96 Ps.97}{true}\chronologyannotationreach{Ps.95.3}{Ps.67 Ps.95 Ps.96 Ps.97}{true}\chronologyannotationreach{Ps.95.4}{Ps.67 Ps.95 Ps.96 Ps.97}{true}\chronologyannotationreach{Ps.95.5}{Ps.67 Ps.95 Ps.96 Ps.97}{true}\chronologyannotationreach{Ps.95.7}{Ps.67 Ps.95 Ps.96 Ps.97}{true}\chronologyannotationreach{Ps.95.8}{Ps.67 Ps.95 Ps.96 Ps.97}{true}\chronologyannotationreach{Ps.95.9}{Ps.67 Ps.95 Ps.96 Ps.97}{true}\chronologyannotationclaim{life-of-christ.nativity}{prophetic-referent}{catholic-traditional-v1}{disputed}{the year of Rome 750 which he styles 3 B.C.}{The Nativity of Our Lord, B.C. 3}; \chronologyannotationreach{Ps.95.1}{Ps.67 Ps.95 Ps.96 Ps.97}{true}\chronologyannotationreach{Ps.95.10}{Ps.67 Ps.95 Ps.96 Ps.97}{true}\chronologyannotationreach{Ps.95.3}{Ps.67 Ps.95 Ps.96 Ps.97}{true}\chronologyannotationreach{Ps.95.4}{Ps.67 Ps.95 Ps.96 Ps.97}{true}\chronologyannotationreach{Ps.95.5}{Ps.67 Ps.95 Ps.96 Ps.97}{true}\chronologyannotationreach{Ps.95.7}{Ps.67 Ps.95 Ps.96 Ps.97}{true}\chronologyannotationreach{Ps.95.8}{Ps.67 Ps.95 Ps.96 Ps.97}{true}\chronologyannotationreach{Ps.95.9}{Ps.67 Ps.95 Ps.96 Ps.97}{true}\chronologyannotationclaim{life-of-christ.nativity}{prophetic-referent}{catholic-traditional-v1}{disputed}{probably the year 7 B.C.}{The Nativity of Our Lord, c. B.C. 7}.}%
}
% second-reading: Second reading
\expandafter\def\csname triptychchronologyannotation@second-reading\endcsname{%
  \chronologyannotationgroup{composition}{disputed}{\textbf{Composition} -- disputed: \chronologyannotationreach{1Thess.1.1}{1Thess}{true}\chronologyannotationreach{1Thess.1.2}{1Thess}{true}\chronologyannotationreach{1Thess.1.3}{1Thess}{true}\chronologyannotationreach{1Thess.1.4}{1Thess}{true}\chronologyannotationreach{1Thess.1.5}{1Thess}{true}\chronologyannotationclaim{composition.first-epistle-to-the-thessalonians}{composition}{catholic-traditional-v1}{disputed}{(1 and 2 Thessalonians), 52}{A.D. 52}; \chronologyannotationreach{1Thess.1.1}{1Thess}{true}\chronologyannotationreach{1Thess.1.2}{1Thess}{true}\chronologyannotationreach{1Thess.1.3}{1Thess}{true}\chronologyannotationreach{1Thess.1.4}{1Thess}{true}\chronologyannotationreach{1Thess.1.5}{1Thess}{true}\chronologyannotationclaim{composition.first-epistle-to-the-thessalonians}{composition}{catholic-traditional-v1}{disputed}{in the year 53 or 54 according to the commonly received scheme of Pauline chronology}{A.D. 53--54}.}%
}
% acclamation: Gospel acclamation
\expandafter\def\csname triptychchronologyannotation@acclamation\endcsname{%
  \chronologyannotationgroup{composition}{disputed}{\textbf{Composition} -- disputed: \chronologyannotationreach{Phil.2.15}{Phil}{true}\chronologyannotationreach{Phil.2.16}{Phil}{true}\chronologyannotationclaim{composition.epistle-to-the-philippians}{composition}{catholic-traditional-v1}{disputed}{(Philemon; Colossians; Ephesians; Philippians), 61}{A.D. 61}; \chronologyannotationreach{Phil.2.15}{Phil}{true}\chronologyannotationreach{Phil.2.16}{Phil}{true}\chronologyannotationclaim{composition.epistle-to-the-philippians}{composition}{catholic-traditional-v1}{disputed}{at Rome (A.D. 62-64)}{A.D. 62--64}.}%
}
% gospel: Gospel
\expandafter\def\csname triptychchronologyannotation@gospel\endcsname{%
  \chronologyannotationgroup{narrated-event}{research-pending}{\textbf{Event} -- No narrated-event date in the chronology corpus.}%
  \space \chronologyannotationgroup{composition}{disputed}{\textbf{Composition} -- disputed: \chronologyannotationreach{Matt.22.15}{Matt}{true}\chronologyannotationreach{Matt.22.16}{Matt}{true}\chronologyannotationreach{Matt.22.17}{Matt}{true}\chronologyannotationreach{Matt.22.18}{Matt}{true}\chronologyannotationreach{Matt.22.19}{Matt}{true}\chronologyannotationreach{Matt.22.20}{Matt}{true}\chronologyannotationreach{Matt.22.21}{Matt}{true}\chronologyannotationclaim{composition.gospel-of-matthew}{composition}{catholic-traditional-v1}{disputed}{about A.D. 38-45}{c. A.D. 38--45}; \chronologyannotationreach{Matt.22.15}{Matt}{true}\chronologyannotationreach{Matt.22.16}{Matt}{true}\chronologyannotationreach{Matt.22.17}{Matt}{true}\chronologyannotationreach{Matt.22.18}{Matt}{true}\chronologyannotationreach{Matt.22.19}{Matt}{true}\chronologyannotationreach{Matt.22.20}{Matt}{true}\chronologyannotationreach{Matt.22.21}{Matt}{true}\chronologyannotationclaim{composition.gospel-of-matthew}{composition}{catholic-traditional-v1}{disputed}{about the year 40-42}{c. A.D. 40--42}; \chronologyannotationreach{Matt.22.15}{Matt}{true}\chronologyannotationreach{Matt.22.16}{Matt}{true}\chronologyannotationreach{Matt.22.17}{Matt}{true}\chronologyannotationreach{Matt.22.18}{Matt}{true}\chronologyannotationreach{Matt.22.19}{Matt}{true}\chronologyannotationreach{Matt.22.20}{Matt}{true}\chronologyannotationreach{Matt.22.21}{Matt}{true}\chronologyannotationclaim{composition.gospel-of-matthew}{composition}{catholic-traditional-v1}{disputed}{the years 40-45}{A.D. 40--45}; \chronologyannotationreach{Matt.22.15}{Matt}{true}\chronologyannotationreach{Matt.22.16}{Matt}{true}\chronologyannotationreach{Matt.22.17}{Matt}{true}\chronologyannotationreach{Matt.22.18}{Matt}{true}\chronologyannotationreach{Matt.22.19}{Matt}{true}\chronologyannotationreach{Matt.22.20}{Matt}{true}\chronologyannotationreach{Matt.22.21}{Matt}{true}\chronologyannotationclaim{composition.gospel-of-matthew}{composition}{catholic-traditional-v1}{disputed}{about the year 60-68}{c. A.D. 60--68}; \chronologyannotationreach{Matt.22.15}{Matt}{true}\chronologyannotationreach{Matt.22.16}{Matt}{true}\chronologyannotationreach{Matt.22.17}{Matt}{true}\chronologyannotationreach{Matt.22.18}{Matt}{true}\chronologyannotationreach{Matt.22.19}{Matt}{true}\chronologyannotationreach{Matt.22.20}{Matt}{true}\chronologyannotationreach{Matt.22.21}{Matt}{true}\chronologyannotationclaim{composition.gospel-of-matthew}{composition}{catholic-traditional-v1}{disputed}{about the years 64-67}{c. A.D. 64--67}; \chronologyannotationreach{Matt.22.15}{Matt}{true}\chronologyannotationreach{Matt.22.16}{Matt}{true}\chronologyannotationreach{Matt.22.17}{Matt}{true}\chronologyannotationreach{Matt.22.18}{Matt}{true}\chronologyannotationreach{Matt.22.19}{Matt}{true}\chronologyannotationreach{Matt.22.20}{Matt}{true}\chronologyannotationreach{Matt.22.21}{Matt}{true}\chronologyannotationclaim{composition.gospel-of-matthew}{composition}{catholic-traditional-v1}{disputed}{about the year 50}{c. A.D. 50}.}%
}
% communion-psalm: Communion: psalm option
\expandafter\def\csname triptychchronologyannotation@communion-psalm\endcsname{%
  \chronologyannotationgroup{traditional-attribution}{disputed}{\textbf{Traditional attribution} -- disputed: \chronologyannotationreach{Ps.32.18}{Ps.32}{true}\chronologyannotationreach{Ps.32.19}{Ps.32}{true}\chronologyannotationclaim{israel.monarchy.david-traditional-era}{traditional-attribution}{catholic-traditional-v1}{disputed}{reigned from 1055 to 1015 B.C.}{David (reign in the usual chronology), B.C. 1055--1015}.}%
  \space \chronologyannotationgroup{historical-setting}{research-pending}{\textbf{Historical setting} -- Occasion date unresolved.}%
  \space \chronologyannotationgroup{composition}{preferred}{\textbf{Composition}: \chronologyannotationreach{Ps.32.18}{Ps}{true}\chronologyannotationreach{Ps.32.19}{Ps}{true}\chronologyannotationclaim{critical.psalms.latest-composition-boundary}{composition}{catholic-critical-v1}{preferred}{before the Maccabean period, around 165 B.C.; no individual psalm can be dated securely}{Before c. 165 B.C.}}%
}
% communion-mark: Communion: Mark option
\expandafter\def\csname triptychchronologyannotation@communion-mark\endcsname{%
  \chronologyannotationgroup{composition}{disputed}{\textbf{Composition} -- disputed: \chronologyannotationreach{Mark.10.45}{Mark}{true}\chronologyannotationclaim{composition.gospel-of-mark}{composition}{catholic-traditional-v1}{disputed}{the tenth or twelfth year after the Ascension (A.D. 38-40)}{A.D. 38--40}; \chronologyannotationreach{Mark.10.45}{Mark}{true}\chronologyannotationclaim{composition.gospel-of-mark}{composition}{catholic-traditional-v1}{disputed}{A.D. 40}{A.D. 40}; \chronologyannotationreach{Mark.10.45}{Mark}{true}\chronologyannotationclaim{composition.gospel-of-mark}{composition}{catholic-traditional-v1}{disputed}{the third year of Claudius (A.D. 43)}{A.D. 43}; \chronologyannotationreach{Mark.10.45}{Mark}{true}\chronologyannotationclaim{composition.gospel-of-mark}{composition}{catholic-traditional-v1}{disputed}{somewhere between A.D. 50 and 67}{A.D. 50--67}; \chronologyannotationreach{Mark.10.45}{Mark}{true}\chronologyannotationclaim{composition.gospel-of-mark}{composition}{catholic-traditional-v1}{disputed}{between 60 and 70}{A.D. 60--70}; \chronologyannotationreach{Mark.10.45}{Mark}{true}\chronologyannotationclaim{composition.gospel-of-mark}{composition}{catholic-traditional-v1}{disputed}{after Matthew, and not as late as the destruction of Jerusalem}{After Matthew, and not as late as the destruction of Jerusalem}.}%
}
% presentation: orientation-v3
\newcommand{\chronologyorientationrow}[2]{%
  \ifcsname triptychchronologyorientation@#1\endcsname
    \csname triptychchronologyorientation@#1\endcsname{#2}%
  \else
    \PackageError{triptych}{No chronology orientation row for #1}{}%
  \fi}
\expandafter\def\csname triptychchronologyorientation@row-8e2bbfdf31f6907d\endcsname#1{%
Entrance antiphon: Cf. Ps. 16:6,8 (Vg.; Heb. 17) & Historical setting & Occasion date unresolved\\*[.1em]%
}
\expandafter\def\csname triptychchronologyorientation@row-a1ffcff124ad6db6\endcsname#1{%
 & Composition: Psalms & \chronologyannotationreach{Ps.16.6}{Ps}{true}\chronologyannotationreach{Ps.16.8}{Ps}{true}\chronologyannotationclaim{critical.psalms.latest-composition-boundary}{composition}{catholic-critical-v1}{preferred}{before the Maccabean period, around 165 B.C.; no individual psalm can be dated securely}{Before c.~165 B.C.}\\[.3em]%
}
\expandafter\def\csname triptychchronologyorientation@row-f7c093da7b8aceaa\endcsname#1{%
Communion: psalm option: Cf. Ps. 32:18--19 (Vg.; Heb. 33) & Traditional attribution: David's reign; Hebron, then Jerusalem & Disputed: \chronologyannotationreach{Ps.32.18}{Ps.32}{true}\chronologyannotationreach{Ps.32.19}{Ps.32}{true}\chronologyannotationclaim{israel.monarchy.david-traditional-era}{traditional-attribution}{catholic-traditional-v1}{disputed}{reigned from 1055 to 1015 B.C.}{B.C.~1055--1015}\\*[.1em]%
}
\expandafter\def\csname triptychchronologyorientation@row-27504bebae67aab8\endcsname#1{%
 & Historical setting & Occasion date unresolved\\*[.1em]%
}
\expandafter\def\csname triptychchronologyorientation@row-3507714e01df9670\endcsname#1{%
 & Composition: Psalms & \chronologyannotationreach{Ps.32.18}{Ps}{true}\chronologyannotationreach{Ps.32.19}{Ps}{true}\chronologyannotationclaim{critical.psalms.latest-composition-boundary}{composition}{catholic-critical-v1}{preferred}{before the Maccabean period, around 165 B.C.; no individual psalm can be dated securely}{Before c.~165 B.C.}\\[.3em]%
}
\expandafter\def\csname triptychchronologyorientation@row-2b7bf42bbda519bd\endcsname#1{%
Responsorial Psalm: Ps. 96:1,3,4--5,7--8,9,10a,10c (Heb.; Vg. 95), Ps. 96:7b (Heb.; Vg. 95) & Superscription setting: The raising of the second Temple by Zorobabel; Temple at Jerusalem & \chronologyannotationreach{Ps.95.1}{Ps.95}{true}\chronologyannotationreach{Ps.95.10}{Ps.95}{true}\chronologyannotationreach{Ps.95.3}{Ps.95}{true}\chronologyannotationreach{Ps.95.4}{Ps.95}{true}\chronologyannotationreach{Ps.95.5}{Ps.95}{true}\chronologyannotationreach{Ps.95.7}{Ps.95}{true}\chronologyannotationreach{Ps.95.8}{Ps.95}{true}\chronologyannotationreach{Ps.95.9}{Ps.95}{true}\chronologyannotationclaim{israel.restoration.temple-raised-by-zorobabel}{superscription-setting}{catholic-traditional-v1}{preferred}{(537 B.C.)}{B.C.~537}\\*[.1em]%
}
\expandafter\def\csname triptychchronologyorientation@row-c79d335b1500eeee\endcsname#1{%
 & Composition: Psalms & \chronologyannotationreach{Ps.95.1}{Ps}{true}\chronologyannotationreach{Ps.95.10}{Ps}{true}\chronologyannotationreach{Ps.95.3}{Ps}{true}\chronologyannotationreach{Ps.95.4}{Ps}{true}\chronologyannotationreach{Ps.95.5}{Ps}{true}\chronologyannotationreach{Ps.95.7}{Ps}{true}\chronologyannotationreach{Ps.95.8}{Ps}{true}\chronologyannotationreach{Ps.95.9}{Ps}{true}\chronologyannotationclaim{critical.psalms.latest-composition-boundary}{composition}{catholic-critical-v1}{preferred}{before the Maccabean period, around 165 B.C.; no individual psalm can be dated securely}{Before c.~165 B.C.}\\*[.1em]%
}
\expandafter\def\csname triptychchronologyorientation@row-271393e4d5eba396\endcsname#1{%
 & Prophetic referent: The Nativity of Our Lord; Bethlehem as the Christian referent & Disputed: \chronologyannotationreach{Ps.95.1}{Ps.67 Ps.95 Ps.96 Ps.97}{true}\chronologyannotationreach{Ps.95.10}{Ps.67 Ps.95 Ps.96 Ps.97}{true}\chronologyannotationreach{Ps.95.3}{Ps.67 Ps.95 Ps.96 Ps.97}{true}\chronologyannotationreach{Ps.95.4}{Ps.67 Ps.95 Ps.96 Ps.97}{true}\chronologyannotationreach{Ps.95.5}{Ps.67 Ps.95 Ps.96 Ps.97}{true}\chronologyannotationreach{Ps.95.7}{Ps.67 Ps.95 Ps.96 Ps.97}{true}\chronologyannotationreach{Ps.95.8}{Ps.67 Ps.95 Ps.96 Ps.97}{true}\chronologyannotationreach{Ps.95.9}{Ps.67 Ps.95 Ps.96 Ps.97}{true}\chronologyannotationclaim{life-of-christ.nativity}{prophetic-referent}{catholic-traditional-v1}{disputed}{the year of Rome 750 which he styles 3 B.C.}{B.C.~3}; \chronologyannotationreach{Ps.95.1}{Ps.67 Ps.95 Ps.96 Ps.97}{true}\chronologyannotationreach{Ps.95.10}{Ps.67 Ps.95 Ps.96 Ps.97}{true}\chronologyannotationreach{Ps.95.3}{Ps.67 Ps.95 Ps.96 Ps.97}{true}\chronologyannotationreach{Ps.95.4}{Ps.67 Ps.95 Ps.96 Ps.97}{true}\chronologyannotationreach{Ps.95.5}{Ps.67 Ps.95 Ps.96 Ps.97}{true}\chronologyannotationreach{Ps.95.7}{Ps.67 Ps.95 Ps.96 Ps.97}{true}\chronologyannotationreach{Ps.95.8}{Ps.67 Ps.95 Ps.96 Ps.97}{true}\chronologyannotationreach{Ps.95.9}{Ps.67 Ps.95 Ps.96 Ps.97}{true}\chronologyannotationclaim{life-of-christ.nativity}{prophetic-referent}{catholic-traditional-v1}{disputed}{probably the year 7 B.C.}{c.~B.C.~7}\\[.3em]%
}
\expandafter\def\csname triptychchronologyorientation@row-5bd350c42030e13d\endcsname#1{%
First reading: Is. 45:1,4--6 & Traditional attribution: Isaias (ministry in Souvay's traditional account); Judah and Jerusalem & \chronologyannotationreach{Is.45.1}{Is}{true}\chronologyannotationreach{Is.45.4}{Is}{true}\chronologyannotationreach{Is.45.5}{Is}{true}\chronologyannotationreach{Is.45.6}{Is}{true}\chronologyannotationclaim{israel.prophets.isaias-traditional-era}{traditional-attribution}{catholic-traditional-v1}{preferred}{from the closing year of Ozias, King of Juda (740 B.C.); no prophecies appear to be later than 701 B.C.}{B.C.~740--701}\\*[.1em]%
}
\expandafter\def\csname triptychchronologyorientation@row-61813b60283fb6d8\endcsname#1{%
 & Prophecy given: The prophetic activity behind Isaiah 40-55 in the NABRE introduction; exiles from Babylon; restoration to Jerusalem & \chronologyannotationreach{Is.45.1}{Is.40 Is.41 Is.42 Is.43 Is.44 Is.45 Is.46 Is.47 Is.48 Is.49 Is.50 Is.51 Is.52 Is.53 Is.54 Is.55}{true}\chronologyannotationreach{Is.45.4}{Is.40 Is.41 Is.42 Is.43 Is.44 Is.45 Is.46 Is.47 Is.48 Is.49 Is.50 Is.51 Is.52 Is.53 Is.54 Is.55}{true}\chronologyannotationreach{Is.45.5}{Is.40 Is.41 Is.42 Is.43 Is.44 Is.45 Is.46 Is.47 Is.48 Is.49 Is.50 Is.51 Is.52 Is.53 Is.54 Is.55}{true}\chronologyannotationreach{Is.45.6}{Is.40 Is.41 Is.42 Is.43 Is.44 Is.45 Is.46 Is.47 Is.48 Is.49 Is.50 Is.51 Is.52 Is.53 Is.54 Is.55}{true}\chronologyannotationclaim{israel.exile.second-isaias-oracles}{prophecy-given}{catholic-critical-v1}{preferred}{toward the end of the Babylonian exile}{Toward the end of the Babylonian exile}\\[.3em]%
}
\expandafter\def\csname triptychchronologyorientation@row-839bf079c07a920a\endcsname#1{%
Gospel: Matt. 22:15--21 & Narrated event: Jesus answers about tribute to Caesar in the Temple at Jerusalem & Date unresolved\\*[.1em]%
}
\expandafter\def\csname triptychchronologyorientation@row-bf53325f9e4771a0\endcsname#1{%
 & Composition: Matthew; Semitic original associated with Palestine & Disputed: \chronologyannotationreach{Matt.22.15}{Matt}{true}\chronologyannotationreach{Matt.22.16}{Matt}{true}\chronologyannotationreach{Matt.22.17}{Matt}{true}\chronologyannotationreach{Matt.22.18}{Matt}{true}\chronologyannotationreach{Matt.22.19}{Matt}{true}\chronologyannotationreach{Matt.22.20}{Matt}{true}\chronologyannotationreach{Matt.22.21}{Matt}{true}\chronologyannotationclaim{composition.gospel-of-matthew}{composition}{catholic-traditional-v1}{disputed}{about A.D. 38-45}{c.~A.D.~38--45}; \chronologyannotationreach{Matt.22.15}{Matt}{true}\chronologyannotationreach{Matt.22.16}{Matt}{true}\chronologyannotationreach{Matt.22.17}{Matt}{true}\chronologyannotationreach{Matt.22.18}{Matt}{true}\chronologyannotationreach{Matt.22.19}{Matt}{true}\chronologyannotationreach{Matt.22.20}{Matt}{true}\chronologyannotationreach{Matt.22.21}{Matt}{true}\chronologyannotationclaim{composition.gospel-of-matthew}{composition}{catholic-traditional-v1}{disputed}{about the year 40-42}{c.~A.D.~40--42}; \chronologyannotationreach{Matt.22.15}{Matt}{true}\chronologyannotationreach{Matt.22.16}{Matt}{true}\chronologyannotationreach{Matt.22.17}{Matt}{true}\chronologyannotationreach{Matt.22.18}{Matt}{true}\chronologyannotationreach{Matt.22.19}{Matt}{true}\chronologyannotationreach{Matt.22.20}{Matt}{true}\chronologyannotationreach{Matt.22.21}{Matt}{true}\chronologyannotationclaim{composition.gospel-of-matthew}{composition}{catholic-traditional-v1}{disputed}{the years 40-45}{A.D.~40--45}; \chronologyannotationreach{Matt.22.15}{Matt}{true}\chronologyannotationreach{Matt.22.16}{Matt}{true}\chronologyannotationreach{Matt.22.17}{Matt}{true}\chronologyannotationreach{Matt.22.18}{Matt}{true}\chronologyannotationreach{Matt.22.19}{Matt}{true}\chronologyannotationreach{Matt.22.20}{Matt}{true}\chronologyannotationreach{Matt.22.21}{Matt}{true}\chronologyannotationclaim{composition.gospel-of-matthew}{composition}{catholic-traditional-v1}{disputed}{about the year 60-68}{c.~A.D.~60--68}; \chronologyannotationreach{Matt.22.15}{Matt}{true}\chronologyannotationreach{Matt.22.16}{Matt}{true}\chronologyannotationreach{Matt.22.17}{Matt}{true}\chronologyannotationreach{Matt.22.18}{Matt}{true}\chronologyannotationreach{Matt.22.19}{Matt}{true}\chronologyannotationreach{Matt.22.20}{Matt}{true}\chronologyannotationreach{Matt.22.21}{Matt}{true}\chronologyannotationclaim{composition.gospel-of-matthew}{composition}{catholic-traditional-v1}{disputed}{about the years 64-67}{c.~A.D.~64--67}; \chronologyannotationreach{Matt.22.15}{Matt}{true}\chronologyannotationreach{Matt.22.16}{Matt}{true}\chronologyannotationreach{Matt.22.17}{Matt}{true}\chronologyannotationreach{Matt.22.18}{Matt}{true}\chronologyannotationreach{Matt.22.19}{Matt}{true}\chronologyannotationreach{Matt.22.20}{Matt}{true}\chronologyannotationreach{Matt.22.21}{Matt}{true}\chronologyannotationclaim{composition.gospel-of-matthew}{composition}{catholic-traditional-v1}{disputed}{about the year 50}{c.~A.D.~50}\\[.3em]%
}
\expandafter\def\csname triptychchronologyorientation@row-afd0d6a0b63b1029\endcsname#1{%
Communion: Mark option: Mark 10:45 & Composition: Mark; Rome traditionally; Egypt reported by Chrysostom & Disputed: \chronologyannotationreach{Mark.10.45}{Mark}{true}\chronologyannotationclaim{composition.gospel-of-mark}{composition}{catholic-traditional-v1}{disputed}{the tenth or twelfth year after the Ascension (A.D. 38-40)}{A.D.~38--40}; \chronologyannotationreach{Mark.10.45}{Mark}{true}\chronologyannotationclaim{composition.gospel-of-mark}{composition}{catholic-traditional-v1}{disputed}{A.D. 40}{A.D.~40}; \chronologyannotationreach{Mark.10.45}{Mark}{true}\chronologyannotationclaim{composition.gospel-of-mark}{composition}{catholic-traditional-v1}{disputed}{the third year of Claudius (A.D. 43)}{A.D.~43}; \chronologyannotationreach{Mark.10.45}{Mark}{true}\chronologyannotationclaim{composition.gospel-of-mark}{composition}{catholic-traditional-v1}{disputed}{somewhere between A.D. 50 and 67}{A.D.~50--67}; \chronologyannotationreach{Mark.10.45}{Mark}{true}\chronologyannotationclaim{composition.gospel-of-mark}{composition}{catholic-traditional-v1}{disputed}{between 60 and 70}{A.D.~60--70}; \chronologyannotationreach{Mark.10.45}{Mark}{true}\chronologyannotationclaim{composition.gospel-of-mark}{composition}{catholic-traditional-v1}{disputed}{after Matthew, and not as late as the destruction of Jerusalem}{After Matthew, and not as late as the destruction of Jerusalem}\\[.3em]%
}
\expandafter\def\csname triptychchronologyorientation@row-dd97b891425fc821\endcsname#1{%
Gospel acclamation: Phil. 2:15d,16a & Composition: Philippians; Rome to Philippi in the cited account & Disputed: \chronologyannotationreach{Phil.2.15}{Phil}{true}\chronologyannotationreach{Phil.2.16}{Phil}{true}\chronologyannotationclaim{composition.epistle-to-the-philippians}{composition}{catholic-traditional-v1}{disputed}{(Philemon; Colossians; Ephesians; Philippians), 61}{A.D.~61}; \chronologyannotationreach{Phil.2.15}{Phil}{true}\chronologyannotationreach{Phil.2.16}{Phil}{true}\chronologyannotationclaim{composition.epistle-to-the-philippians}{composition}{catholic-traditional-v1}{disputed}{at Rome (A.D. 62-64)}{A.D.~62--64}\\[.3em]%
}
\expandafter\def\csname triptychchronologyorientation@row-71c06b197ea8f727\endcsname#1{%
Second reading: 1 Thess. 1:1--5b & Composition: First Thessalonians; Corinth to Thessalonica & Disputed: \chronologyannotationreach{1Thess.1.1}{1Thess}{true}\chronologyannotationreach{1Thess.1.2}{1Thess}{true}\chronologyannotationreach{1Thess.1.3}{1Thess}{true}\chronologyannotationreach{1Thess.1.4}{1Thess}{true}\chronologyannotationreach{1Thess.1.5}{1Thess}{true}\chronologyannotationclaim{composition.first-epistle-to-the-thessalonians}{composition}{catholic-traditional-v1}{disputed}{(1 and 2 Thessalonians), 52}{A.D.~52}; \chronologyannotationreach{1Thess.1.1}{1Thess}{true}\chronologyannotationreach{1Thess.1.2}{1Thess}{true}\chronologyannotationreach{1Thess.1.3}{1Thess}{true}\chronologyannotationreach{1Thess.1.4}{1Thess}{true}\chronologyannotationreach{1Thess.1.5}{1Thess}{true}\chronologyannotationclaim{composition.first-epistle-to-the-thessalonians}{composition}{catholic-traditional-v1}{disputed}{in the year 53 or 54 according to the commonly received scheme of Pauline chronology}{A.D.~53--54}\\[.3em]%
}

% Generated. Do not edit.
% Complete orientation-v3 table; wording, dates and order belong to the corpus.
\begin{chronologysharedtable}
\chronologyorientationrow{row-8e2bbfdf31f6907d}{}
\chronologyorientationrow{row-a1ffcff124ad6db6}{}
\midrule
\chronologyorientationrow{row-f7c093da7b8aceaa}{}
\chronologyorientationrow{row-27504bebae67aab8}{}
\chronologyorientationrow{row-3507714e01df9670}{}
\midrule
\chronologyorientationrow{row-2b7bf42bbda519bd}{}
\chronologyorientationrow{row-c79d335b1500eeee}{}
\chronologyorientationrow{row-271393e4d5eba396}{}
\midrule
\chronologyorientationrow{row-5bd350c42030e13d}{}
\chronologyorientationrow{row-61813b60283fb6d8}{}
\midrule
\chronologyorientationrow{row-839bf079c07a920a}{}
\chronologyorientationrow{row-bf53325f9e4771a0}{}
\midrule
\chronologyorientationrow{row-afd0d6a0b63b1029}{}
\midrule
\chronologyorientationrow{row-dd97b891425fc821}{}
\midrule
\chronologyorientationrow{row-71c06b197ea8f727}{}
\end{chronologysharedtable}
{\footnotesize Psalm numbering: Vg. = Vulgate; Heb. = Hebrew.\par}


\subsection{Evidence for each date}
\label{sec:date-location-evidence}
{\small For each passage, as the Scriptural Date and Location table names it: what is dated, how the corpus weighs each answer, the source's own words, and the sources cited. A subject the corpus does not date is reported as such and given no date.\par}
{\small
\par\smallskip\noindent\textit{Entrance antiphon: Cf. Ps. 16:6,8 (Vg.; Heb. 17)}\par\nobreak
\begin{itemize}[leftmargin=1.2em,itemsep=1pt,topsep=1pt,parsep=0pt]
\item \textbf{Historical setting}. No date: the chronology corpus has not yet dated this (research pending).
\item \textbf{Composition} (preferred). Psalms: Before c.~165 B.C. (``before the Maccabean period, around 165 B.C.; no individual psalm can be dated securely''; Catholic critical chronology; New American Bible, Revised Edition --- USCCB web state (2026), Psalms introduction).
\end{itemize}
\par\smallskip\noindent\textit{Communion: psalm option: Cf. Ps. 32:18--19 (Vg.; Heb. 33)}\par\nobreak
\begin{itemize}[leftmargin=1.2em,itemsep=1pt,topsep=1pt,parsep=0pt]
\item \textbf{Traditional attribution} (disputed). David's reign in the usual chronology: B.C.~1055--1015 (Corbett also reports writers dating the reign 30--50 years later) (``reigned from 1055 to 1015 B.C.''; The Catholic Encyclopedia, Volume 4 (New York, 1908), article 'David, King,' printed p. 642, opening paragraph). Place: Hebron, then Jerusalem (The Holy Bible, Douay-Rheims (Challoner revision), Project Gutenberg eBook 1581: 2 Kings 5:5).
\item \textbf{Historical setting}. As under Entrance antiphon: Cf. Ps. 16:6,8 (Vg.; Heb. 17).
\item \textbf{Composition}. As under Entrance antiphon: Cf. Ps. 16:6,8 (Vg.; Heb. 17).
\end{itemize}
\par\smallskip\noindent\textit{Responsorial Psalm: Ps. 96:1,3,4--5,7--8,9,10a,10c (Heb.; Vg. 95), Ps. 96:7b (Heb.; Vg. 95)}\par\nobreak
\begin{itemize}[leftmargin=1.2em,itemsep=1pt,topsep=1pt,parsep=0pt]
\item \textbf{Superscription setting} (preferred). The raising of the second Temple by Zorobabel: B.C.~537 (``(537 B.C.)''; The Catholic Encyclopedia, Volume 14 (New York, 1912)). Place: Temple at Jerusalem (The Holy Bible, Douay-Rheims (Challoner revision), Project Gutenberg eBook 1581: Ps. 95:1, 1 Esd. 3:8).
\item \textbf{Composition}. As under Entrance antiphon: Cf. Ps. 16:6,8 (Vg.; Heb. 17).
\item \textbf{Prophetic referent} (disputed). The Nativity of Our Lord: B.C.~3 (``the year of Rome 750 which he styles 3 B.C.''; The Catholic Encyclopedia, Volume 3 (New York, 1908)); c.~B.C.~7 (``probably the year 7 B.C.''; The Catholic Encyclopedia, Volume 3 (New York, 1908)). Place: Bethlehem as the Christian referent (The Holy Bible, Douay-Rheims (Challoner revision), Project Gutenberg eBook 1581: Matt. 2:1).
\end{itemize}
\par\smallskip\noindent\textit{First reading: Is. 45:1,4--6}\par\nobreak
\begin{itemize}[leftmargin=1.2em,itemsep=1pt,topsep=1pt,parsep=0pt]
\item \textbf{Traditional attribution} (preferred). Isaias (ministry in Souvay's traditional account): B.C.~740--701 (``from the closing year of Ozias, King of Juda (740 B.C.); no prophecies appear to be later than 701 B.C.''; The Catholic Encyclopedia, Volume 8 (New York, 1910)). Place: Judah and Jerusalem (The Holy Bible, Douay-Rheims (Challoner revision), Project Gutenberg eBook 1581: Is. 1:1).
\item \textbf{Prophecy given} (preferred). The prophetic activity behind Isaiah 40-55 in the NABRE introduction: Toward the end of the Babylonian exile (Catholic critical chronology; New American Bible, Revised Edition: USCCB web state (2026), Isaiah introduction, especially the final two prose paragraphs before the outline). Place: exiles from Babylon; restoration to Jerusalem (The Holy Bible, Douay-Rheims (Challoner revision), Project Gutenberg eBook 1581: Is. 48:20, 52:9).
\end{itemize}
\par\smallskip\noindent\textit{Gospel: Matt. 22:15--21}\par\nobreak
\begin{itemize}[leftmargin=1.2em,itemsep=1pt,topsep=1pt,parsep=0pt]
\item \textbf{Narrated event}. Jesus answers about tribute to Caesar in the Temple at Jerusalem (The Holy Bible, Douay-Rheims (Challoner revision), Project Gutenberg eBook 1581: Matt. 21:10, 21:23, 22:15, 22:17, 22:21). No date: the chronology corpus has not yet dated this (research pending).
\item \textbf{Composition} (disputed). The Gospel of St Matthew: c.~A.D.~38--45 (``about A.D. 38-45''; The Catholic Encyclopedia, Volume 10 (1911); The Catholic Encyclopedia, Volume 14 (New York, 1912)); c.~A.D.~40--42 (``about the year 40-42''; The Catholic Encyclopedia, Volume 10 (1911)); A.D.~40--45 (``the years 40-45''; The Catholic Encyclopedia, Volume 10 (1911)); c.~A.D.~60--68 (``about the year 60-68''; The Catholic Encyclopedia, Volume 10 (1911)); c.~A.D.~64--67 (``about the years 64-67''; The Catholic Encyclopedia, Volume 10 (1911)); c.~A.D.~50 (``about the year 50''; The Catholic Encyclopedia, Volume 14 (New York, 1912)). Place: Semitic original associated with Palestine (The Catholic Encyclopedia, Volume 10 (1911)).
\end{itemize}
\par\smallskip\noindent\textit{Communion: Mark option: Mark 10:45}\par\nobreak
\begin{itemize}[leftmargin=1.2em,itemsep=1pt,topsep=1pt,parsep=0pt]
\item \textbf{Composition} (disputed). The Gospel of St Mark: A.D.~38--40 (``the tenth or twelfth year after the Ascension (A.D. 38-40)''; The Catholic Encyclopedia, Volume 9 (1910)); A.D.~40 (The Catholic Encyclopedia, Volume 9 (1910)); A.D.~43 (``the third year of Claudius (A.D. 43)''; The Catholic Encyclopedia, Volume 9 (1910)); A.D.~50--67 (``somewhere between A.D. 50 and 67''; The Catholic Encyclopedia, Volume 9 (1910)); A.D.~60--70 (``between 60 and 70''; The Catholic Encyclopedia, Volume 14 (New York, 1912)); After Matthew, and not as late as the destruction of Jerusalem (De auctore, de tempore compositionis et de historica veritate Evangeliorum secundum Marcum et secundum Lucam (1912), responsa-i-ix). Place: Rome traditionally; Egypt reported by Chrysostom (The Catholic Encyclopedia, Volume 9 (1910)).
\end{itemize}
\par\smallskip\noindent\textit{Gospel acclamation: Phil. 2:15d,16a}\par\nobreak
\begin{itemize}[leftmargin=1.2em,itemsep=1pt,topsep=1pt,parsep=0pt]
\item \textbf{Composition} (disputed). The Epistle to the Philippians: A.D.~61 (``(Philemon; Colossians; Ephesians; Philippians), 61''; The Catholic Encyclopedia, Volume 11 (New York, 1911)); A.D.~62--64 (``at Rome (A.D. 62-64)''; The Catholic Encyclopedia, Volume 12 (1911)). Place: Rome to Philippi in the cited account (The Catholic Encyclopedia, Volume 12 (1911); The Holy Bible, Douay-Rheims (Challoner revision), Project Gutenberg eBook 1581: Phil. 1:1).
\end{itemize}
\par\smallskip\noindent\textit{Second reading: 1 Thess. 1:1--5b}\par\nobreak
\begin{itemize}[leftmargin=1.2em,itemsep=1pt,topsep=1pt,parsep=0pt]
\item \textbf{Composition} (disputed). The First Epistle to the Thessalonians: A.D.~52 (``(1 and 2 Thessalonians), 52''; The Catholic Encyclopedia, Volume 11 (New York, 1911)); A.D.~53--54 (``in the year 53 or 54 according to the commonly received scheme of Pauline chronology''; The Catholic Encyclopedia, Volume 14 (New York, 1912)). Place: Corinth to Thessalonica (The Catholic Encyclopedia, Volume 14 (New York, 1912); The Holy Bible, Douay-Rheims (Challoner revision), Project Gutenberg eBook 1581: 1 Thess. 1:1).
\end{itemize}
\par}
